\section{Matchertext and its security relevance}
\label{sec:relevance}

The narrow waist of \cref{sec:bg:waist} is a claim about \emph{embedding},
not about security.
Matchertext was designed so that a string in one language can be placed in another
verbatim, without escaping.
Its stated benefit is that the writer of embedded text need reason about only one
language's rules at a time~\cite{matchertext}.
That is a usability argument.
This section states the discipline precisely enough to reason about.
It explains why a property invented for embedding bears on injection at all, and
sets out the questions the rest of the paper answers.

\subsection{The discipline}
\label{sec:relevance:discipline}

Matchertext is a purely syntactic rule over an alphabet $\Sigma$ and a set $\Pi$
of open/close character pairs: the \emph{matchers} must match.
Write $O$ and $C$ for the openers and closers of $\Pi$, and \emph{nonmatchers}
for every other character.
The language $L$ of matchertext strings is defined
inductively: any string of nonmatchers is in $L$, and if $m_1,m_2,m_3\in L$ and
$(o,c)\in\Pi$ then $m_1\,o\,m_2\,c\,m_3\in L$.
The concrete configuration takes $\Sigma$ to be Unicode and $\Pi$ to be the three
ASCII pairs \verb|()|, \verb|[]|, and \verb|{}|.
The angle brackets are excluded because they are used unmatched as inequality
operators far too often to constrain~\cite{matchertext}.
Nothing else is fixed: the discipline assigns no meaning to any character and is
deliberately oblivious to the syntax of any particular language.

Three results carry over from the original development, and the rest of this paper
uses only these.
First, \emph{closure}, also called always-embeddability: if $m\in L$ and
$(s_o,s_c)$ is a matched delimiter pair, then $s_o\,m\,s_c\in L$.
A valid matchertext string placed inside matched delimiters therefore yields a
valid matchertext string, at any nesting depth and with no escaping or expansion.
Second, membership in $L$ is decided by a single linear scan, the \textsc{Verify}
recognizer.
It maintains a stack of unclosed openers and accepts exactly when every closer
matches its opener and the stack ends empty.
Third, a total encoder \textsc{ToMatchertext} maps \emph{any} string into $L$ by
escaping only its unmatched matchers.
Free-form input that legitimately contains an unmatched bracket therefore need
not be rejected.
Closure and the correctness of \textsc{Verify} are machine-checked; so is
\textsc{ToMatchertext}'s landing in $L$ on every input
(\cref{sec:appendix:tomatchertext}), leaving only its decoder round-trip informal.
For the rationale behind $\Pi$, the hosting syntaxes, and the empirical
compliance study, see~\cite{matchertext}.

\subsection{From embedding to security}
\label{sec:relevance:security}

Verbatim embedding and injection resistance are not the same property, and the
paper's claim turns on the difference between them.

Injection, as \cref{sec:motivation} describes it, is a boundary failure.
A value meant to occupy one syntactic region escapes it, because a delimiter
inside the value collides with the delimiter that should contain it.
Conventional defenses attack the \emph{collision}: they transform the value so
that its dangerous characters can no longer be confused with the host's.
That is why they must know the host's syntax, and why there is a different one
for every host and context.
Matchertext attacks the \emph{boundary} instead.
Suppose a host delimits an untrusted value with a matcher pair.
It finds the end of that value by counting matchers, not by scanning for a
terminator.
Then escaping the region requires emitting an unmatched matcher, and by closure
no valid matchertext string contains one.
The characters that drive the classic breakouts, the quote in SQL and the angle
bracket in HTML, are nonmatchers.
The discipline leaves them entirely unconstrained, and they become ordinary data
inside a matched region.
Closure keeps that region balanced whatever nonmatchers the value carries.

Two consequences follow, and they are the reason this is a security question
rather than a syntactic one.
The check is \emph{host-independent}: \textsc{Verify} inspects only the six
matcher characters and knows nothing about SQL, HTML, or LDAP.
The same recognizer therefore serves every host, which is exactly what a
sanitizer cannot do.
And the check is a \emph{verification} rather than a transformation, so its failure
mode is benign.
A rejected value is data that does not embed, whereas a missed case in an escaping
routine \emph{is} the injection.
Together these are the narrow waist applied to security: one uniform obligation on
the value, in place of one bespoke transformation per host-context pair.

\subsection{What must hold, and what we ask}
\label{sec:relevance:questions}

Closure constrains \emph{strings}, whereas security is a claim about
\emph{parsers}, and the second does not follow from the first without assumptions
about the host.
\Cref{sec:threat} states those assumptions; three questions organize the rest of
the paper.

\paragraph{Q1: what is actually guaranteed, and under what assumptions?}
Delimiting a value is not by itself a security property, and ``cannot break out''
is not a specification.
\Cref{sec:threat} states the guarantee as an \emph{inertness} invariant on the
host's parse.
It separates that guarantee from the weaker closure property it is easily
conflated with.
And it reduces the guarantee to a small set of named premises about the host
parser, the interpreter, and the pipeline that reaches them.

\paragraph{Q2: can a real language adopt this, and at what cost?}
The guarantee assumes a matchertext-aware host, and no mainstream language is one
today.
\Cref{sec:resist} develops the argument for the syntaxes where injection actually
occurs.
It then asks what adopting it costs: which hosting options deliver the property,
and what each adds to the trusted computing base.
A final question is how the result compares with the defenses already deployed
for these hosts.

\paragraph{Q3: how much of the real injection landscape could this reach?}
Even granting the property and its adoption, the discipline addresses only
injections that are structural breakouts from matcher-delimitable contexts.
That is a proper subset of what the term covers.
\Cref{sec:corpus,sec:results,sec:prevent} measure the subset against a corpus of
recorded vulnerabilities.
\Cref{sec:prevent} reports what fraction of
attacks with a recoverable payload would be contained.
It also reports what fraction of the landscape is out of reach in principle.
